\section*{TÓM TẮT ĐỒ ÁN\\TỐI ƯU HÓA BẢO MẬT ĐA LỚP VÀ CÂN BẰNG TẢI DỰA TRÊN NGƯỠNG TRONG MÔ HÌNH CAMPUS 3 LỚP}

\begin{itemize}
    \item \textbf{Vấn đề nghiên cứu:}
    
    Đồ án tập trung nghiên cứu, thiết kế và mô phỏng một hạ tầng mạng Campus 3 lớp (Core, Distribution, Access) có tích hợp vùng phi quân sự (DMZ). Mục tiêu cốt lõi là giải quyết bài toán cấp phát dịch vụ nội bộ ra bên ngoài một cách an toàn thông qua các kỹ thuật dịch địa chỉ (NAT/PAT), kiểm soát luồng (ACL) và tự động hóa điều phối lưu lượng (Load Balancing). Báo cáo được cấu trúc gồm 5 chương:
    
        \begin{itemize}
        \item[--] Chương 1: Tổng quan và Thiết kế Topology.
        \item[--] Chương 2: Cấu hình Định tuyến OSPF và NAT/PAT.
        \item[--] Chương 3: Triển khai Bảo mật đa lớp.
        \item[--] Chương 4: Thực nghiệm Cân bằng tải theo ngưỡng.
        \item[--] Chương 5: Phân tích kết quả và Kết luận.
        \end{itemize}
        
    \item \textbf{Các hướng tiếp cận:}
    
        \begin{itemize}
        \item[--] \textit{Lý thuyết:} Nghiên cứu kiến trúc Campus 3-tier, nguyên lý hoạt động của NAT/PAT, cấu trúc ACL và Firewall trên Linux (iptables), giao thức định tuyến OSPF.
        
        \item[--] \textit{Thực hành:} Sử dụng trình giả lập Mininet kết hợp với FRRouting để triển khai mạng diện rộng. Áp dụng Script Python để tự động hóa giám sát băng thông, định tuyến lại (Traffic Steering) và đo lường hiệu năng mạng.
        \end{itemize}
        
    \item \textbf{Cách giải quyết vấn đề:}
    
    Xây dựng một cấu trúc mạng mô phỏng hoàn chỉnh gồm Firewall, Core Router, Distribution và Access Switch. Cấu hình OSPF để đồng bộ bảng định tuyến nội bộ. Triển khai các quy tắc iptables khắt khe để bảo vệ vùng DMZ và cho phép truy cập ra Internet. Viết kịch bản Python giám sát ngưỡng 80\% băng thông để tự động cấu hình lại DNAT, chuyển hướng luồng dữ liệu sang máy chủ dự phòng.
    
    \item \textbf{Kết quả đạt được:}
    
    Hệ thống đã hội tụ định tuyến thành công và thực hiện chuyển đổi địa chỉ (NAT) chính xác (được xác nhận qua Traceroute và bảng NAT). Kịch bản cân bằng tải hoạt động ổn định và trực quan. Các kịch bản đo lường tự động đã cung cấp số liệu thực tế về thông lượng (Throughput), độ trễ (Delay), qua đó chứng minh tác động overhead nhỏ bé của Firewall so với lợi ích bảo mật mang lại.
    
    Đồ án giúp sinh viên nắm vững kỹ năng thiết kế, cấu hình Tường lửa đa lớp, và ứng dụng tự động hóa vào giám sát hệ thống mạng chuyên nghiệp trên nền tảng giả lập Linux.
\end{itemize}

\newpage